% Môn: Toán
% Lớp: 10
% Chương: Chương II. Bất phương trình và hệ bất phương trình bậc nhất hai ẩn
% Bài: Bài 4: Hệ bất phương trình bậc nhất hai ẩn
% Loại: Lý thuyết
% File riêng, không trộn vào de.tex
% Định nghĩa lệnh Lý thuyết — dùng khi biên dịch lt.tex bằng TeX.
% App web đọc lt.tex trực tiếp (bỏ \input file này).
% Trong lt.tex, từ thư mục bài:
%   % Định nghĩa lệnh Lý thuyết — dùng khi biên dịch lt.tex bằng TeX.
% App web đọc lt.tex trực tiếp (bỏ \input file này).
% Trong lt.tex, từ thư mục bài:
%   % Định nghĩa lệnh Lý thuyết — dùng khi biên dịch lt.tex bằng TeX.
% App web đọc lt.tex trực tiếp (bỏ \input file này).
% Trong lt.tex, từ thư mục bài:
%   \input{../../../../_lenh/lythuyet.tex}
% từ gốc repo:
%   \input{ngan-hang/_lenh/lythuyet.tex}
\makeatletter
\providecommand{\indam}[1]{\textbf{#1}}
\providecommand{\chuy}[1]{\par\smallskip\noindent\textit{#1}\par}
\providecommand{\luuy}[1]{\par\smallskip\noindent\textbf{Lưu ý.} #1\par}
\providecommand{\ghichu}[1]{\par\smallskip\noindent\textbf{Ghi chú.} #1\par}
\providecommand{\vidu}[1]{\par\smallskip\noindent\textbf{Ví dụ.} #1\par}
\providecommand{\Blythuyet}{\subsection*{Lý thuyết}}
% Hai cột: kiến thức | hình
\providecommand{\haicotchay}[2]{%
  \par\medskip\noindent
  \begin{minipage}[t]{0.52\linewidth}#1\end{minipage}\hfill
  \begin{minipage}[t]{0.46\linewidth}#2\end{minipage}%
  \par\medskip
}
\providecommand{\kienthuccot}[2]{\haicotchay{#1}{#2}}
% Dự phòng nếu chưa nạp graphicx / caption (không ghi đè bản thật)
\providecommand{\resizebox}[3]{#3}
\providecommand{\captionof}[2]{\par\smallskip\noindent\textit{#2}\par}
\newenvironment{hoatdong}[1]{%
  \par\medskip\noindent\textbf{Hoạt động.} #1\par\smallskip
}{\par\medskip}
\newenvironment{emcobiet}{%
  \par\medskip\noindent\textbf{Em có biết}\par\smallskip
}{\par\medskip}
\newenvironment{emdahoc}{%
  \par\medskip\noindent\textbf{Em đã học}\par\smallskip
}{\par\medskip}
\newenvironment{emcothe}{%
  \par\medskip\noindent\textbf{Em có thể}\par\smallskip
}{\par\medskip}
\newenvironment{kienthuc}{%
  \par\medskip\noindent\textbf{Kiến thức cốt lõi}\par\smallskip
}{\par\medskip}
\newenvironment{traloi}{%
  \par\smallskip\noindent\textit{Trả lời / ghi chú của em}\par
  \vspace{1.2em}%
}{\par\medskip}
\newenvironment{phuongphap}{%
  \par\medskip\noindent\textbf{Phương pháp giải}\par\smallskip
}{\par\medskip}
\newenvironment{vidumau}{%
  \par\medskip\noindent\textbf{Bài mẫu}\par\smallskip
}{\par\medskip}
\newenvironment{dangmau}[1]{%
  \par\medskip\noindent\textbf{Dạng mẫu.} #1\par\smallskip
}{\par\medskip}
\makeatother

% từ gốc repo:
%   % Định nghĩa lệnh Lý thuyết — dùng khi biên dịch lt.tex bằng TeX.
% App web đọc lt.tex trực tiếp (bỏ \input file này).
% Trong lt.tex, từ thư mục bài:
%   \input{../../../../_lenh/lythuyet.tex}
% từ gốc repo:
%   \input{ngan-hang/_lenh/lythuyet.tex}
\makeatletter
\providecommand{\indam}[1]{\textbf{#1}}
\providecommand{\chuy}[1]{\par\smallskip\noindent\textit{#1}\par}
\providecommand{\luuy}[1]{\par\smallskip\noindent\textbf{Lưu ý.} #1\par}
\providecommand{\ghichu}[1]{\par\smallskip\noindent\textbf{Ghi chú.} #1\par}
\providecommand{\vidu}[1]{\par\smallskip\noindent\textbf{Ví dụ.} #1\par}
\providecommand{\Blythuyet}{\subsection*{Lý thuyết}}
% Hai cột: kiến thức | hình
\providecommand{\haicotchay}[2]{%
  \par\medskip\noindent
  \begin{minipage}[t]{0.52\linewidth}#1\end{minipage}\hfill
  \begin{minipage}[t]{0.46\linewidth}#2\end{minipage}%
  \par\medskip
}
\providecommand{\kienthuccot}[2]{\haicotchay{#1}{#2}}
% Dự phòng nếu chưa nạp graphicx / caption (không ghi đè bản thật)
\providecommand{\resizebox}[3]{#3}
\providecommand{\captionof}[2]{\par\smallskip\noindent\textit{#2}\par}
\newenvironment{hoatdong}[1]{%
  \par\medskip\noindent\textbf{Hoạt động.} #1\par\smallskip
}{\par\medskip}
\newenvironment{emcobiet}{%
  \par\medskip\noindent\textbf{Em có biết}\par\smallskip
}{\par\medskip}
\newenvironment{emdahoc}{%
  \par\medskip\noindent\textbf{Em đã học}\par\smallskip
}{\par\medskip}
\newenvironment{emcothe}{%
  \par\medskip\noindent\textbf{Em có thể}\par\smallskip
}{\par\medskip}
\newenvironment{kienthuc}{%
  \par\medskip\noindent\textbf{Kiến thức cốt lõi}\par\smallskip
}{\par\medskip}
\newenvironment{traloi}{%
  \par\smallskip\noindent\textit{Trả lời / ghi chú của em}\par
  \vspace{1.2em}%
}{\par\medskip}
\newenvironment{phuongphap}{%
  \par\medskip\noindent\textbf{Phương pháp giải}\par\smallskip
}{\par\medskip}
\newenvironment{vidumau}{%
  \par\medskip\noindent\textbf{Bài mẫu}\par\smallskip
}{\par\medskip}
\newenvironment{dangmau}[1]{%
  \par\medskip\noindent\textbf{Dạng mẫu.} #1\par\smallskip
}{\par\medskip}
\makeatother

\makeatletter
\providecommand{\indam}[1]{\textbf{#1}}
\providecommand{\chuy}[1]{\par\smallskip\noindent\textit{#1}\par}
\providecommand{\luuy}[1]{\par\smallskip\noindent\textbf{Lưu ý.} #1\par}
\providecommand{\ghichu}[1]{\par\smallskip\noindent\textbf{Ghi chú.} #1\par}
\providecommand{\vidu}[1]{\par\smallskip\noindent\textbf{Ví dụ.} #1\par}
\providecommand{\Blythuyet}{\subsection*{Lý thuyết}}
% Hai cột: kiến thức | hình
\providecommand{\haicotchay}[2]{%
  \par\medskip\noindent
  \begin{minipage}[t]{0.52\linewidth}#1\end{minipage}\hfill
  \begin{minipage}[t]{0.46\linewidth}#2\end{minipage}%
  \par\medskip
}
\providecommand{\kienthuccot}[2]{\haicotchay{#1}{#2}}
% Dự phòng nếu chưa nạp graphicx / caption (không ghi đè bản thật)
\providecommand{\resizebox}[3]{#3}
\providecommand{\captionof}[2]{\par\smallskip\noindent\textit{#2}\par}
\newenvironment{hoatdong}[1]{%
  \par\medskip\noindent\textbf{Hoạt động.} #1\par\smallskip
}{\par\medskip}
\newenvironment{emcobiet}{%
  \par\medskip\noindent\textbf{Em có biết}\par\smallskip
}{\par\medskip}
\newenvironment{emdahoc}{%
  \par\medskip\noindent\textbf{Em đã học}\par\smallskip
}{\par\medskip}
\newenvironment{emcothe}{%
  \par\medskip\noindent\textbf{Em có thể}\par\smallskip
}{\par\medskip}
\newenvironment{kienthuc}{%
  \par\medskip\noindent\textbf{Kiến thức cốt lõi}\par\smallskip
}{\par\medskip}
\newenvironment{traloi}{%
  \par\smallskip\noindent\textit{Trả lời / ghi chú của em}\par
  \vspace{1.2em}%
}{\par\medskip}
\newenvironment{phuongphap}{%
  \par\medskip\noindent\textbf{Phương pháp giải}\par\smallskip
}{\par\medskip}
\newenvironment{vidumau}{%
  \par\medskip\noindent\textbf{Bài mẫu}\par\smallskip
}{\par\medskip}
\newenvironment{dangmau}[1]{%
  \par\medskip\noindent\textbf{Dạng mẫu.} #1\par\smallskip
}{\par\medskip}
\makeatother

% từ gốc repo:
%   % Định nghĩa lệnh Lý thuyết — dùng khi biên dịch lt.tex bằng TeX.
% App web đọc lt.tex trực tiếp (bỏ \input file này).
% Trong lt.tex, từ thư mục bài:
%   % Định nghĩa lệnh Lý thuyết — dùng khi biên dịch lt.tex bằng TeX.
% App web đọc lt.tex trực tiếp (bỏ \input file này).
% Trong lt.tex, từ thư mục bài:
%   \input{../../../../_lenh/lythuyet.tex}
% từ gốc repo:
%   \input{ngan-hang/_lenh/lythuyet.tex}
\makeatletter
\providecommand{\indam}[1]{\textbf{#1}}
\providecommand{\chuy}[1]{\par\smallskip\noindent\textit{#1}\par}
\providecommand{\luuy}[1]{\par\smallskip\noindent\textbf{Lưu ý.} #1\par}
\providecommand{\ghichu}[1]{\par\smallskip\noindent\textbf{Ghi chú.} #1\par}
\providecommand{\vidu}[1]{\par\smallskip\noindent\textbf{Ví dụ.} #1\par}
\providecommand{\Blythuyet}{\subsection*{Lý thuyết}}
% Hai cột: kiến thức | hình
\providecommand{\haicotchay}[2]{%
  \par\medskip\noindent
  \begin{minipage}[t]{0.52\linewidth}#1\end{minipage}\hfill
  \begin{minipage}[t]{0.46\linewidth}#2\end{minipage}%
  \par\medskip
}
\providecommand{\kienthuccot}[2]{\haicotchay{#1}{#2}}
% Dự phòng nếu chưa nạp graphicx / caption (không ghi đè bản thật)
\providecommand{\resizebox}[3]{#3}
\providecommand{\captionof}[2]{\par\smallskip\noindent\textit{#2}\par}
\newenvironment{hoatdong}[1]{%
  \par\medskip\noindent\textbf{Hoạt động.} #1\par\smallskip
}{\par\medskip}
\newenvironment{emcobiet}{%
  \par\medskip\noindent\textbf{Em có biết}\par\smallskip
}{\par\medskip}
\newenvironment{emdahoc}{%
  \par\medskip\noindent\textbf{Em đã học}\par\smallskip
}{\par\medskip}
\newenvironment{emcothe}{%
  \par\medskip\noindent\textbf{Em có thể}\par\smallskip
}{\par\medskip}
\newenvironment{kienthuc}{%
  \par\medskip\noindent\textbf{Kiến thức cốt lõi}\par\smallskip
}{\par\medskip}
\newenvironment{traloi}{%
  \par\smallskip\noindent\textit{Trả lời / ghi chú của em}\par
  \vspace{1.2em}%
}{\par\medskip}
\newenvironment{phuongphap}{%
  \par\medskip\noindent\textbf{Phương pháp giải}\par\smallskip
}{\par\medskip}
\newenvironment{vidumau}{%
  \par\medskip\noindent\textbf{Bài mẫu}\par\smallskip
}{\par\medskip}
\newenvironment{dangmau}[1]{%
  \par\medskip\noindent\textbf{Dạng mẫu.} #1\par\smallskip
}{\par\medskip}
\makeatother

% từ gốc repo:
%   % Định nghĩa lệnh Lý thuyết — dùng khi biên dịch lt.tex bằng TeX.
% App web đọc lt.tex trực tiếp (bỏ \input file này).
% Trong lt.tex, từ thư mục bài:
%   \input{../../../../_lenh/lythuyet.tex}
% từ gốc repo:
%   \input{ngan-hang/_lenh/lythuyet.tex}
\makeatletter
\providecommand{\indam}[1]{\textbf{#1}}
\providecommand{\chuy}[1]{\par\smallskip\noindent\textit{#1}\par}
\providecommand{\luuy}[1]{\par\smallskip\noindent\textbf{Lưu ý.} #1\par}
\providecommand{\ghichu}[1]{\par\smallskip\noindent\textbf{Ghi chú.} #1\par}
\providecommand{\vidu}[1]{\par\smallskip\noindent\textbf{Ví dụ.} #1\par}
\providecommand{\Blythuyet}{\subsection*{Lý thuyết}}
% Hai cột: kiến thức | hình
\providecommand{\haicotchay}[2]{%
  \par\medskip\noindent
  \begin{minipage}[t]{0.52\linewidth}#1\end{minipage}\hfill
  \begin{minipage}[t]{0.46\linewidth}#2\end{minipage}%
  \par\medskip
}
\providecommand{\kienthuccot}[2]{\haicotchay{#1}{#2}}
% Dự phòng nếu chưa nạp graphicx / caption (không ghi đè bản thật)
\providecommand{\resizebox}[3]{#3}
\providecommand{\captionof}[2]{\par\smallskip\noindent\textit{#2}\par}
\newenvironment{hoatdong}[1]{%
  \par\medskip\noindent\textbf{Hoạt động.} #1\par\smallskip
}{\par\medskip}
\newenvironment{emcobiet}{%
  \par\medskip\noindent\textbf{Em có biết}\par\smallskip
}{\par\medskip}
\newenvironment{emdahoc}{%
  \par\medskip\noindent\textbf{Em đã học}\par\smallskip
}{\par\medskip}
\newenvironment{emcothe}{%
  \par\medskip\noindent\textbf{Em có thể}\par\smallskip
}{\par\medskip}
\newenvironment{kienthuc}{%
  \par\medskip\noindent\textbf{Kiến thức cốt lõi}\par\smallskip
}{\par\medskip}
\newenvironment{traloi}{%
  \par\smallskip\noindent\textit{Trả lời / ghi chú của em}\par
  \vspace{1.2em}%
}{\par\medskip}
\newenvironment{phuongphap}{%
  \par\medskip\noindent\textbf{Phương pháp giải}\par\smallskip
}{\par\medskip}
\newenvironment{vidumau}{%
  \par\medskip\noindent\textbf{Bài mẫu}\par\smallskip
}{\par\medskip}
\newenvironment{dangmau}[1]{%
  \par\medskip\noindent\textbf{Dạng mẫu.} #1\par\smallskip
}{\par\medskip}
\makeatother

\makeatletter
\providecommand{\indam}[1]{\textbf{#1}}
\providecommand{\chuy}[1]{\par\smallskip\noindent\textit{#1}\par}
\providecommand{\luuy}[1]{\par\smallskip\noindent\textbf{Lưu ý.} #1\par}
\providecommand{\ghichu}[1]{\par\smallskip\noindent\textbf{Ghi chú.} #1\par}
\providecommand{\vidu}[1]{\par\smallskip\noindent\textbf{Ví dụ.} #1\par}
\providecommand{\Blythuyet}{\subsection*{Lý thuyết}}
% Hai cột: kiến thức | hình
\providecommand{\haicotchay}[2]{%
  \par\medskip\noindent
  \begin{minipage}[t]{0.52\linewidth}#1\end{minipage}\hfill
  \begin{minipage}[t]{0.46\linewidth}#2\end{minipage}%
  \par\medskip
}
\providecommand{\kienthuccot}[2]{\haicotchay{#1}{#2}}
% Dự phòng nếu chưa nạp graphicx / caption (không ghi đè bản thật)
\providecommand{\resizebox}[3]{#3}
\providecommand{\captionof}[2]{\par\smallskip\noindent\textit{#2}\par}
\newenvironment{hoatdong}[1]{%
  \par\medskip\noindent\textbf{Hoạt động.} #1\par\smallskip
}{\par\medskip}
\newenvironment{emcobiet}{%
  \par\medskip\noindent\textbf{Em có biết}\par\smallskip
}{\par\medskip}
\newenvironment{emdahoc}{%
  \par\medskip\noindent\textbf{Em đã học}\par\smallskip
}{\par\medskip}
\newenvironment{emcothe}{%
  \par\medskip\noindent\textbf{Em có thể}\par\smallskip
}{\par\medskip}
\newenvironment{kienthuc}{%
  \par\medskip\noindent\textbf{Kiến thức cốt lõi}\par\smallskip
}{\par\medskip}
\newenvironment{traloi}{%
  \par\smallskip\noindent\textit{Trả lời / ghi chú của em}\par
  \vspace{1.2em}%
}{\par\medskip}
\newenvironment{phuongphap}{%
  \par\medskip\noindent\textbf{Phương pháp giải}\par\smallskip
}{\par\medskip}
\newenvironment{vidumau}{%
  \par\medskip\noindent\textbf{Bài mẫu}\par\smallskip
}{\par\medskip}
\newenvironment{dangmau}[1]{%
  \par\medskip\noindent\textbf{Dạng mẫu.} #1\par\smallskip
}{\par\medskip}
\makeatother

\makeatletter
\providecommand{\indam}[1]{\textbf{#1}}
\providecommand{\chuy}[1]{\par\smallskip\noindent\textit{#1}\par}
\providecommand{\luuy}[1]{\par\smallskip\noindent\textbf{Lưu ý.} #1\par}
\providecommand{\ghichu}[1]{\par\smallskip\noindent\textbf{Ghi chú.} #1\par}
\providecommand{\vidu}[1]{\par\smallskip\noindent\textbf{Ví dụ.} #1\par}
\providecommand{\Blythuyet}{\subsection*{Lý thuyết}}
% Hai cột: kiến thức | hình
\providecommand{\haicotchay}[2]{%
  \par\medskip\noindent
  \begin{minipage}[t]{0.52\linewidth}#1\end{minipage}\hfill
  \begin{minipage}[t]{0.46\linewidth}#2\end{minipage}%
  \par\medskip
}
\providecommand{\kienthuccot}[2]{\haicotchay{#1}{#2}}
% Dự phòng nếu chưa nạp graphicx / caption (không ghi đè bản thật)
\providecommand{\resizebox}[3]{#3}
\providecommand{\captionof}[2]{\par\smallskip\noindent\textit{#2}\par}
\newenvironment{hoatdong}[1]{%
  \par\medskip\noindent\textbf{Hoạt động.} #1\par\smallskip
}{\par\medskip}
\newenvironment{emcobiet}{%
  \par\medskip\noindent\textbf{Em có biết}\par\smallskip
}{\par\medskip}
\newenvironment{emdahoc}{%
  \par\medskip\noindent\textbf{Em đã học}\par\smallskip
}{\par\medskip}
\newenvironment{emcothe}{%
  \par\medskip\noindent\textbf{Em có thể}\par\smallskip
}{\par\medskip}
\newenvironment{kienthuc}{%
  \par\medskip\noindent\textbf{Kiến thức cốt lõi}\par\smallskip
}{\par\medskip}
\newenvironment{traloi}{%
  \par\smallskip\noindent\textit{Trả lời / ghi chú của em}\par
  \vspace{1.2em}%
}{\par\medskip}
\newenvironment{phuongphap}{%
  \par\medskip\noindent\textbf{Phương pháp giải}\par\smallskip
}{\par\medskip}
\newenvironment{vidumau}{%
  \par\medskip\noindent\textbf{Bài mẫu}\par\smallskip
}{\par\medskip}
\newenvironment{dangmau}[1]{%
  \par\medskip\noindent\textbf{Dạng mẫu.} #1\par\smallskip
}{\par\medskip}
\makeatother




\subsubsection{Bất phương trình bậc nhất hai ẩn}

\begin{hoatdong}
\chuy{Nhắc lại kiến thức về bất phương trình bậc nhất một ẩn và phương trình bậc nhất hai ẩn.}
\begin{itemize}
    \item Bất phương trình bậc nhất một ẩn có dạng như thế nào? Cho ví dụ.
    \item Phương trình bậc nhất hai ẩn có dạng như thế nào? Cho ví dụ.
\end{itemize}
\chuy{Xét bài toán sau:}
Một cửa hàng bán hai loại áo sơ mi: loại A giá 200 nghìn đồng/chiếc và loại B giá 300 nghìn đồng/chiếc. Cửa hàng muốn bán được ít nhất 1 triệu đồng tiền áo. Gọi $x$ là số áo loại A bán được và $y$ là số áo loại B bán được.
\begin{itemize}
    \item Viết biểu thức thể hiện tổng số tiền bán được.
    \item Viết bất phương trình biểu thị điều kiện cửa hàng muốn bán được ít nhất 1 triệu đồng.
\end{itemize}
\end{hoatdong}

\begin{kienthuc}
\chuy{Định nghĩa}
Bất phương trình bậc nhất hai ẩn $x, y$ có một trong các dạng sau:
\begin{itemize}
    \item $ax + by + c > 0$
    \item $ax + by + c < 0$
    \item $ax + by + c \ge 0$
    \item $ax + by + c \le 0$
\end{itemize}
trong đó $a, b, c$ là các số thực đã cho, với $a$ và $b$ không đồng thời bằng $0$.
\chuy{Nghiệm của bất phương trình bậc nhất hai ẩn}
Cặp số $(x_0; y_0)$ được gọi là một nghiệm của bất phương trình bậc nhất hai ẩn nếu thay $x = x_0$ và $y = y_0$ vào bất phương trình ta được một khẳng định đúng.
\end{kienthuc}

\begin{emdahoc}
\chuy{Kiểm tra cặp số là nghiệm của bất phương trình bậc nhất hai ẩn.}
Để kiểm tra xem một cặp số $(x_0; y_0)$ có phải là nghiệm của bất phương trình $ax + by + c > 0$ (hoặc $<, \ge, \le$) hay không, ta thay $x_0, y_0$ vào bất phương trình và kiểm tra tính đúng sai của khẳng định thu được.
\end{emdahoc}

\begin{emcothe}
\chuy{Ví dụ:}
Cặp số nào sau đây là nghiệm của bất phương trình $x - 6y - 4 > 0$?
\begin{itemize}
    \item A. $(1; 0)$
    \item B. $(10; 1)$
    \item C. $(0; -1)$
    \item D. $(5; 0)$
\end{itemize}
\chuy{Giải:}
\begin{itemize}
    \item Với $(1; 0)$: $1 - 6(0) - 4 = 1 - 4 = -3$. Vì $-3 \ngtr 0$, nên $(1; 0)$ không là nghiệm.
    \item Với $(10; 1)$: $10 - 6(1) - 4 = 10 - 6 - 4 = 0$. Vì $0 \ngtr 0$, nên $(10; 1)$ không là nghiệm.
    \item Với $(0; -1)$: $0 - 6(-1) - 4 = 6 - 4 = 2$. Vì $2 > 0$, nên $(0; -1)$ là nghiệm.
    \item Với $(5; 0)$: $5 - 6(0) - 4 = 5 - 4 = 1$. Vì $1 > 0$, nên $(5; 0)$ là nghiệm.
\end{itemize}
Trong các lựa chọn trên, có hai cặp số là nghiệm. Nếu đây là câu hỏi trắc nghiệm chọn một đáp án đúng duy nhất, cần kiểm tra lại đề bài hoặc các lựa chọn. Với câu hỏi này, cả C và D đều đúng.
\end{emcothe}

\subsubsection{Miền nghiệm của bất phương trình bậc nhất hai ẩn}

\begin{hoatdong}
\chuy{Nhắc lại về đồ thị hàm số bậc nhất $y = ax + b$.}
\begin{itemize}
    \item Đồ thị của hàm số $y = ax + b$ là gì?
    \item Đường thẳng $ax + by + c = 0$ chia mặt phẳng tọa độ thành mấy phần?
\end{itemize}
\chuy{Xét bất phương trình $x + y - 2 > 0$.}
\begin{itemize}
    \item Vẽ đường thẳng $d: x + y - 2 = 0$ trên mặt phẳng tọa độ.
    \item Lấy một điểm $O(0; 0)$ không nằm trên đường thẳng $d$. Thay tọa độ điểm $O$ vào bất phương trình $x + y - 2 > 0$. Khẳng định thu được đúng hay sai?
    \item Lấy một điểm $A(3; 0)$ không nằm trên đường thẳng $d$. Thay tọa độ điểm $A$ vào bất phương trình $x + y - 2 > 0$. Khẳng định thu được đúng hay sai?
\end{itemize}
\end{hoatdong}

\begin{kienthuc}
\chuy{Miền nghiệm của bất phương trình bậc nhất hai ẩn}
Miền nghiệm của bất phương trình bậc nhất hai ẩn $ax + by + c > 0$ (hoặc $<, \ge, \le$) là tập hợp tất cả các điểm $(x; y)$ trên mặt phẳng tọa độ thỏa mãn bất phương trình đó.
\chuy{Cách xác định miền nghiệm}
Để xác định miền nghiệm của bất phương trình bậc nhất hai ẩn $ax + by + c > 0$ (hoặc $<, \ge, \le$), ta thực hiện các bước sau:
\begin{itemize}
    \item \textbf{Bước 1:} Vẽ đường thẳng $d: ax + by + c = 0$. Đường thẳng này chia mặt phẳng tọa độ thành hai nửa mặt phẳng.
    \item \textbf{Bước 2:} Chọn một điểm $M(x_0; y_0)$ không nằm trên đường thẳng $d$ (thường chọn gốc tọa độ $O(0; 0)$ nếu nó không thuộc $d$).
    \item \textbf{Bước 3:} Thay tọa độ điểm $M$ vào bất phương trình.
    \begin{itemize}
        \item Nếu ta được một khẳng định đúng, thì nửa mặt phẳng chứa điểm $M$ là miền nghiệm của bất phương trình.
        \item Nếu ta được một khẳng định sai, thì nửa mặt phẳng không chứa điểm $M$ là miền nghiệm của bất phương trình.
    \end{itemize}
    \item \textbf{Bước 4:} Kết luận.
    \begin{itemize}
        \item Nếu bất phương trình có dấu $>$ hoặc $<$, thì miền nghiệm là nửa mặt phẳng không chứa đường thẳng $d$ (đường thẳng $d$ được vẽ bằng nét đứt).
        \item Nếu bất phương trình có dấu $\ge$ hoặc $\le$, thì miền nghiệm là nửa mặt phẳng chứa đường thẳng $d$ (đường thẳng $d$ được vẽ bằng nét liền).
    \end{itemize}
\end{itemize}
\end{kienthuc}

\begin{emcothe}
\chuy{Ví dụ:}
Miền nghiệm của bất phương trình $-2x + y \le -5$ là nửa mặt phẳng chứa điểm nào trong các điểm sau?
\begin{itemize}
    \item A. $(0; 0)$
    \item B. $(1; 1)$
    \item C. $(3; 0)$
    \item D. $(2; -1)$
\end{itemize}
\chuy{Giải:}
Xét bất phương trình $-2x + y \le -5$.
\begin{itemize}
    \item \textbf{Bước 1:} Vẽ đường thẳng $d: -2x + y = -5$.
    \item \textbf{Bước 2:} Chọn điểm $O(0; 0)$.
    \item \textbf{Bước 3:} Thay tọa độ điểm $O(0; 0)$ vào bất phương trình: $-2(0) + 0 \le -5 \Leftrightarrow 0 \le -5$. Đây là một khẳng định sai.
    \item \textbf{Bước 4:} Vậy miền nghiệm là nửa mặt phẳng không chứa điểm $O(0; 0)$.
\end{itemize}
Bây giờ ta kiểm tra các điểm trong các lựa chọn:
\begin{itemize}
    \item A. $(0; 0)$: Đã kiểm tra, không thuộc miền nghiệm.
    \item B. $(1; 1)$: $-2(1) + 1 = -1$. Vì $-1 \not\le -5$, nên $(1; 1)$ không thuộc miền nghiệm.
    \item C. $(3; 0)$: $-2(3) + 0 = -6$. Vì $-6 \le -5$, nên $(3; 0)$ thuộc miền nghiệm.
    \item D. $(2; -1)$: $-2(2) + (-1) = -4 - 1 = -5$. Vì $-5 \le -5$, nên $(2; -1)$ thuộc miền nghiệm.
\end{itemize}
Trong trường hợp này, cả C và D đều là các điểm thuộc miền nghiệm. Nếu đây là câu hỏi trắc nghiệm chọn một đáp án đúng duy nhất, cần kiểm tra lại đề bài hoặc các lựa chọn.
\end{emcothe}

\subsubsection{Hệ bất phương trình bậc nhất hai ẩn}

\begin{hoatdong}
\chuy{Nhắc lại về hệ phương trình bậc nhất hai ẩn.}
\begin{itemize}
    \item Hệ phương trình bậc nhất hai ẩn có dạng như thế nào?
    \item Nghiệm của hệ phương trình bậc nhất hai ẩn là gì?
\end{itemize}
\chuy{Xét bài toán thực tế:}
Một người nông dân muốn trồng hai loại cây: cà chua và dưa chuột. Diện tích đất trồng cà chua không quá 5 hecta. Diện tích đất trồng dưa chuột không quá 4 hecta. Tổng diện tích đất trồng không quá 7 hecta. Gọi $x$ là diện tích trồng cà chua và $y$ là diện tích trồng dưa chuột.
\begin{itemize}
    \item Viết các bất phương trình biểu thị các điều kiện trên.
    \item Các bất phương trình này tạo thành một hệ như thế nào?
\end{itemize}
\end{hoatdong}

\begin{kienthuc}
\chuy{Định nghĩa}
Hệ bất phương trình bậc nhất hai ẩn là một tập hợp gồm hai hay nhiều bất phương trình bậc nhất hai ẩn.
\chuy{Nghiệm của hệ bất phương trình bậc nhất hai ẩn}
Cặp số $(x_0; y_0)$ được gọi là một nghiệm của hệ bất phương trình bậc nhất hai ẩn nếu nó là nghiệm của tất cả các bất phương trình trong hệ.
\end{kienthuc}

\begin{emdahoc}
\chuy{Kiểm tra cặp số là nghiệm của hệ bất phương trình bậc nhất hai ẩn.}
Để kiểm tra xem một cặp số $(x_0; y_0)$ có phải là nghiệm của một hệ bất phương trình hay không, ta thay $x_0, y_0$ vào từng bất phương trình trong hệ. Nếu tất cả các bất phương trình đều đúng, thì $(x_0; y_0)$ là nghiệm của hệ. Nếu có ít nhất một bất phương trình sai, thì $(x_0; y_0)$ không là nghiệm của hệ.
\end{emdahoc}

\begin{emcothe}
\chuy{Ví dụ:}
Cặp số nào sau đây là nghiệm của hệ bất phương trình $\begin{cases} x + y \ge 2 \\ x - y < 1 \end{cases}$?
\begin{itemize}
    \item A. $(0; 0)$
    \item B. $(1; 1)$
    \item C. $(2; 0)$
    \item D. $(0; 2)$
\end{itemize}
\chuy{Giải:}
\begin{itemize}
    \item A. Với $(0; 0)$:
    $0 + 0 \ge 2 \Leftrightarrow 0 \ge 2$ (Sai). Vậy $(0; 0)$ không là nghiệm của hệ.
    \item B. Với $(1; 1)$:
    $1 + 1 \ge 2 \Leftrightarrow 2 \ge 2$ (Đúng).
    $1 - 1 < 1 \Leftrightarrow 0 < 1$ (Đúng).
    Vậy $(1; 1)$ là nghiệm của hệ.
    \item C. Với $(2; 0)$:
    $2 + 0 \ge 2 \Leftrightarrow 2 \ge 2$ (Đúng).
    $2 - 0 < 1 \Leftrightarrow 2 < 1$ (Sai). Vậy $(2; 0)$ không là nghiệm của hệ.
    \item D. Với $(0; 2)$:
    $0 + 2 \ge 2 \Leftrightarrow 2 \ge 2$ (Đúng).
    $0 - 2 < 1 \Leftrightarrow -2 < 1$ (Đúng).
    Vậy $(0; 2)$ là nghiệm của hệ.
\end{itemize}
Trong các lựa chọn trên, có hai cặp số là nghiệm. Nếu đây là câu hỏi trắc nghiệm chọn một đáp án đúng duy nhất, cần kiểm tra lại đề bài hoặc các lựa chọn.
\end{emcothe}

\subsubsection{Miền nghiệm của hệ bất phương trình bậc nhất hai ẩn}

\begin{kienthuc}
\chuy{Định nghĩa}
Miền nghiệm của một hệ bất phương trình bậc nhất hai ẩn là giao của các miền nghiệm của tất cả các bất phương trình trong hệ.
\chuy{Cách xác định miền nghiệm của hệ bất phương trình}
Để xác định miền nghiệm của hệ bất phương trình bậc nhất hai ẩn, ta thực hiện các bước sau:
\begin{itemize}
    \item \textbf{Bước 1:} Xác định miền nghiệm của từng bất phương trình trong hệ.
    \item \textbf{Bước 2:} Giao các miền nghiệm đó lại. Phần giao nhau (phần chung) chính là miền nghiệm của hệ bất phương trình.
\end{itemize}
\end{kienthuc}

\begin{emcothe}
\chuy{Ví dụ:}
Xác định miền nghiệm của hệ bất phương trình $\begin{cases} x \ge 0 \\ y \ge 0 \\ x + y \le 4 \end{cases}$.
\chuy{Giải:}
\begin{itemize}
    \item Miền nghiệm của $x \ge 0$ là nửa mặt phẳng bên phải trục tung (bao gồm trục tung).
    \item Miền nghiệm của $y \ge 0$ là nửa mặt phẳng phía trên trục hoành (bao gồm trục hoành).
    \item Miền nghiệm của $x + y \le 4$:
    \begin{itemize}
        \item Vẽ đường thẳng $d: x + y = 4$.
        \item Chọn điểm $O(0; 0)$. Thay vào bất phương trình: $0 + 0 \le 4 \Leftrightarrow 0 \le 4$ (Đúng).
        \item Vậy miền nghiệm là nửa mặt phẳng chứa gốc tọa độ $O(0; 0)$ (bao gồm đường thẳng $d$).
    \end{itemize}
\end{itemize}
Giao của ba miền nghiệm này là tam giác OAB với $O(0;0)$, $A(4;0)$ và $B(0;4)$ (bao gồm cả các cạnh và miền trong tam giác).
\end{emcothe}

\subsubsection{Ứng dụng của hệ bất phương trình bậc nhất hai ẩn trong bài toán thực tế}

\begin{kienthuc}
\chuy{Lập bất phương trình/hệ bất phương trình từ bài toán thực tế}
Để lập bất phương trình hoặc hệ bất phương trình từ bài toán thực tế, ta cần:
\begin{itemize}
    \item \textbf{Bước 1:} Đọc kĩ đề bài, xác định các đại lượng cần tìm và đặt biến cho chúng (thường là $x, y$).
    \item \textbf{Bước 2:} Dựa vào các điều kiện, ràng buộc trong bài toán để thiết lập các bất phương trình tương ứng. Lưu ý các từ khóa như "không quá", "ít nhất", "tối đa", "tối thiểu", "không vượt quá", "không nhỏ hơn",...
    \item \textbf{Bước 3:} Kết hợp các bất phương trình thành một hệ bất phương trình.
\end{itemize}
\end{kienthuc}

\begin{emcothe}
\chuy{Ví dụ:}
Bạn Mai có 33 nghìn đồng để đi mua vở. Vở loại I có giá 5000 đồng một cuốn, vở loại II có giá 7000 đồng một cuốn. Tính số quyển vở nhiều nhất bạn Mai có thể mua sao cho bạn có cả hai loại vở.
\chuy{Giải:}
\begin{itemize}
    \item Gọi $x$ là số cuốn vở loại I bạn Mai mua ($x \in \mathbb{N}^*$).
    \item Gọi $y$ là số cuốn vở loại II bạn Mai mua ($y \in \mathbb{N}^*$).
    \item Tổng số tiền mua vở là $5000x + 7000y$.
    \item Vì bạn Mai có 33 nghìn đồng, nên tổng số tiền mua vở không vượt quá 33000 đồng: $5000x + 7000y \le 33000$.
    \item Chia cả hai vế cho 1000, ta được bất phương trình: $5x + 7y \le 33$.
    \item Vì bạn Mai phải có cả hai loại vở, nên $x \ge 1$ và $y \ge 1$.
\end{itemize}
Ta có hệ bất phương trình:
$\begin{cases}
    5x + 7y \le 33 \\
    x \ge 1 \\
    y \ge 1 \\
    x, y \in \mathbb{N}
\end{cases}$
Ta cần tìm $(x; y)$ thỏa mãn hệ và $x+y$ đạt giá trị lớn nhất.
Thử các giá trị $x, y$ nguyên dương:
\begin{itemize}
    \item Nếu $x=1$: $5(1) + 7y \le 33 \Rightarrow 7y \le 28 \Rightarrow y \le 4$. Các cặp $(1;1), (1;2), (1;3), (1;4)$ thỏa mãn. Tổng số vở: $2, 3, 4, 5$.
    \item Nếu $x=2$: $5(2) + 7y \le 33 \Rightarrow 10 + 7y \le 33 \Rightarrow 7y \le 23 \Rightarrow y \le 3.28...$. Vậy $y$ có thể là $1, 2, 3$. Các cặp $(2;1), (2;2), (2;3)$ thỏa mãn. Tổng số vở: $3, 4, 5$.
    \item Nếu $x=3$: $5(3) + 7y \le 33 \Rightarrow 15 + 7y \le 33 \Rightarrow 7y \le 18 \Rightarrow y \le 2.57...$. Vậy $y$ có thể là $1, 2$. Các cặp $(3;1), (3;2)$ thỏa mãn. Tổng số vở: $4, 5$.
    \item Nếu $x=4$: $5(4) + 7y \le 33 \Rightarrow 20 + 7y \le 33 \Rightarrow 7y \le 13 \Rightarrow y \le 1.85...$. Vậy $y$ có thể là $1$. Cặp $(4;1)$ thỏa mãn. Tổng số vở: $5$.
    \item Nếu $x=5$: $5(5) + 7y \le 33 \Rightarrow 25 + 7y \le 33 \Rightarrow 7y \le 8 \Rightarrow y \le 1.14...$. Vậy $y$ có thể là $1$. Cặp $(5;1)$ thỏa mãn. Tổng số vở: $6$.
    \item Nếu $x=6$: $5(6) + 7y \le 33 \Rightarrow 30 + 7y \le 33 \Rightarrow 7y \le 3 \Rightarrow y \le 0.42...$. Không có $y \ge 1$ thỏa mãn.
\end{itemize}
Các cặp $(x;y)$ thỏa mãn là $(1;1), (1;2), (1;3), (1;4), (2;1), (2;2), (2;3), (3;1), (3;2), (4;1), (5;1)$.
Tổng số vở $x+y$ lớn nhất là $1+4=5$, $2+3=5$, $3+2=5$, $4+1=5$, $5+1=6$.
Vậy số quyển vở nhiều nhất bạn Mai có thể mua là 6 quyển (5 cuốn loại I và 1 cuốn loại II).
\end{emcothe}

\subsubsection{Bài toán tối ưu hóa trên miền nghiệm}
\begin{kienthuc}
\chuy{Bài toán tối ưu hóa}
Trong nhiều bài toán thực tế, sau khi lập được hệ bất phương trình và xác định miền nghiệm, ta cần tìm giá trị lớn nhất (GTLN) hoặc giá trị nhỏ nhất (GTNN) của một biểu thức $F(x, y) = Ax + By$ trên miền nghiệm đó.
\chuy{Phương pháp giải}
\begin{itemize}
    \item \textbf{Bước 1:} Lập hệ bất phương trình từ bài toán thực tế và xác định miền nghiệm $D$.
    \item \textbf{Bước 2:} Xác định các đỉnh (các điểm cực trị) của miền nghiệm $D$.
    \item \textbf{Bước 3:} Tính giá trị của biểu thức $F(x, y)$ tại các đỉnh của miền nghiệm $D$.
    \item \textbf{Bước 4:} So sánh các giá trị thu được để tìm GTLN hoặc GTNN của $F(x, y)$ trên miền nghiệm $D$.
\end{itemize}
\chuy{Lưu ý:} Nếu miền nghiệm là một đa giác lồi (miền đóng và bị chặn), thì GTLN và GTNN của biểu thức $F(x, y) = Ax + By$ sẽ đạt được tại một trong các đỉnh của đa giác đó.
\end{kienthuc}

\begin{emcothe}
\chuy{Ví dụ:}
Cho biểu thức $F=4x - 2y$ với $(x;y)$ là các điểm thuộc miền trong hình tứ giác như hình vẽ. Tìm khẳng định đúng.
\begin{center}
 \begin{tikzpicture}[scale=1]
    % Vẽ lưới
    \draw [gray!50,thin] (-1,-1) grid (5,4);
    % Vẽ các điểm
    \filldraw (0,0) circle (1.5pt) node[below left] {$O(0,0)$};
    \filldraw (4,0) circle (1.5pt) node[below right] {$A(4,0)$};
    \filldraw (5,3) circle (1.5pt) node[above right] {$B(5,3)$};
    \filldraw (1,4) circle (1.5pt) node[above left] {$C(1,4)$};
    % Vẽ tứ giác
    \draw[thick, blue, fill=blue!10, opacity=0.7] (0,0) -- (4,0) -- (5,3) -- (1,4) -- cycle;
    % Ghi chú
    \node at (2.5, 2) {Miền nghiệm $D$};
 \end{tikzpicture}
\end{center}
\begin{itemize}
    \item A. Giá trị lớn nhất của $F$ là 16.
    \item B. Giá trị nhỏ nhất của $F$ là -6.
    \item C. Giá trị lớn nhất của $F$ là 14.
    \item D. Giá trị nhỏ nhất của $F$ là -8.
\end{itemize}
\chuy{Giải:}
Miền nghiệm $D$ là tứ giác $OABC$ với các đỉnh là $O(0;0)$, $A(4;0)$, $B(5;3)$, $C(1;4)$.
Ta tính giá trị của biểu thức $F = 4x - 2y$ tại các đỉnh:
\begin{itemize}
    \item Tại $O(0;0)$: $F = 4(0) - 2(0) = 0$.
    \item Tại $A(4;0)$: $F = 4(4) - 2(0) = 16$.
    \item Tại $B(5;3)$: $F = 4(5) - 2(3) = 20 - 6 = 14$.
    \item Tại $C(1;4)$: $F = 4(1) - 2(4) = 4 - 8 = -4$.
\end{itemize}
So sánh các giá trị: $0, 16, 14, -4$.
\begin{itemize}
    \item Giá trị lớn nhất của $F$ là 16 (đạt tại $A(4;0)$).
    \item Giá trị nhỏ nhất của $F$ là -4 (đạt tại $C(1;4)$).
\end{itemize}
Kiểm tra các khẳng định:
\begin{itemize}
    \item A. Giá trị lớn nhất của $F$ là 16. (Đúng)
    \item B. Giá trị nhỏ nhất của $F$ là -6. (Sai, là -4)
    \item C. Giá trị lớn nhất của $F$ là 14. (Sai, là 16)
    \item D. Giá trị nhỏ nhất của $F$ là -8. (Sai, là -4)
\end{itemize}
Vậy khẳng định đúng là A.
\end{emcothe}

\begin{emcobiet}
\chuy{Kiểm tra tính đúng sai của các khẳng định về bất phương trình/hệ bất phương trình bậc nhất hai ẩn và bài toán thực tế.}
Để kiểm tra tính đúng sai của một khẳng định, ta cần áp dụng các định nghĩa và phương pháp đã học:
\begin{itemize}
    \item Đối với khẳng định về nghiệm: Thay cặp số vào bất phương trình/hệ và kiểm tra.
    \item Đối với khẳng định về miền nghiệm: Xác định miền nghiệm và kiểm tra xem điểm/khu vực được nhắc đến có thuộc miền nghiệm hay không.
    \item Đối với khẳng định về bài toán thực tế: Lập hệ bất phương trình, giải và kiểm tra kết quả.
\end{itemize}
\chuy{Ví dụ:}
Xét bất phương trình $-3x + 3y - 3 \ge 0$. Xét tính đúng-sai của các khẳng định sau.
\choiceTF
{ \True Cặp số $(1; 6)$ là một nghiệm của bất phương trình. }
{ \False Cặp số $(-6; -8)$ là nghiệm của bất phương trình. }
\chuy{Giải:}
\begin{itemize}
    \item Với khẳng định 1: Cặp số $(1; 6)$ là một nghiệm của bất phương trình.
    Thay $x=1, y=6$ vào bất phương trình: $-3(1) + 3(6) - 3 = -3 + 18 - 3 = 12$.
    Vì $12 \ge 0$ là khẳng định đúng, nên $(1; 6)$ là nghiệm của bất phương trình. Khẳng định 1 là \textbf{Đúng}.
    \item Với khẳng định 2: Cặp số $(-6; -8)$ là nghiệm của bất phương trình.
    Thay $x=-6, y=-8$ vào bất phương trình: $-3(-6) + 3(-8) - 3 = 18 - 24 - 3 = -9$.
    Vì $-9 \ge 0$ là khẳng định sai, nên $(-6; -8)$ không là nghiệm của bất phương trình. Khẳng định 2 là \textbf{Sai}.
\end{itemize}
\end{emcobiet}
